This element allows simulation of a longitudinal impedance using a
``broad-band'' resonator or an impedance function specified in a file.
The impedance is defined as the Fourier transform of the wake function
\begin{equation}
Z(\omega) = \int_{-\infty}^{+\infty} e^{-i \omega t} W(t) dt
\end{equation}
where $i = \sqrt{-1}$, $W(t)=0$ for $t<0$, and $W(t)$ has units of $V/C$.

For a resonator impedance, the functional form is
\begin{equation}
Z(\omega) = \frac{R_s}{1 + iQ(\frac{\omega}{\omega_r} - \frac{\omega_r}{\omega})},
\end{equation}
where $R_s$ is the shunt impedance in $Ohms$, $Q$ is the quality
factor, and $\omega_r$ is the resonant frequency.

When providing an impedance in a file, the user must be careful to conform to these
conventions. In addition, the units of the frequency column must be Hz, while the units
of the impedance components must be Ohms.
At present, {\tt elegant} does not check the units for correctness.

Other notes:
\begin{enumerate}
\item The frequency data required from the input file is {\em not} $\omega$, but rather
  $f = \omega/(2 \pi)$.
\item The default smoothing setting ({\tt SG\_HALFWIDTH=4}), may apply too much smoothing.
  It is recommended that the user vary this parameter if smoothing is employed.
\item Impedance data can be created from a wake function using the script \verb|wake2impedance|,
  which is supplied with \verb|elegant|. This script also illustrates how to scale the data 
  with the frequency spacing. The script uses \verb|sddsfft|, which produces a
  folded FFT ($f\ge 0$) from a real function. The folded FFT representation involves multiplying the
  non-DC terms by 2. \verb|elegant| expects this and internally multiplies the DC term by 2 as well.
\item Using the broad-brand resonator model can often result in a very large number of bins
 being used, as {\tt elegant} will try to resolve the resonance peak and achieve the desired
 bin spacing. This can result in poor performance, particularly for the parallel version.
\item Wake output is available only in the serial version.
\item By default, the binning grid is centered on the bunch
  ($t_{\rm min} = \langle t\rangle - n_b\,dt/2$) and bin indices are
  computed by truncation ($i_b = \lfloor (t-t_{\rm min})/dt\rfloor$),
  so bin {\em left edges\/} sit at $t_{\rm min} + i_b\,dt$.  Setting
  {\tt BIN\_CENTERED=1} switches to the convention used by {\tt
  ZTRANSVERSE} and {\tt IMPEDANCE}: the bunch is placed one bin in
  from the left edge of the grid, and bin {\em centers\/} sit at
  $t_{\rm min} + i_b\,dt$ (with $i_b$ computed by rounding).  The two
  conventions give physically equivalent results that typically
  differ at the half-bin sampling level ($\sim 10^{-6}$ relative);
  use {\tt BIN\_CENTERED=1} when combining a {\tt ZLONGIT} element
  with a {\tt ZTRANSVERSE} or {\tt IMPEDANCE} element where you want
  the bunch to be binned identically in both.
\end{enumerate}

Bunched-mode application of the impedance is possible using specially-prepared input
beams. 
See Section \ref{sect:bunchedBeams} for details.
The use of bunched mode for any particular \verb|ZLONGIT| element is controlled using the \verb|BUNCHED_BEAM_MODE| parameter.
